% slide-ready summary for the lecture box 'Research from our group' (max 6 lines); generated by the notebook
\begin{itemize}\setlength{\itemsep}{2pt}
  \item \textbf{Prețurile energiei și bursele din UE} (Grecu, Cramer, Pele și Lessmann, 2025, NAJEF): DCC-GARCH și modele Markov; indicele ECE este cel mai slab corelat cu petrolul (max 0,53 în 2020, min $-$0,11 în 2022, Tabelul~2); probabilitatea de persistență a regimurilor depășește 98\,\% (Tabelul~4).
  \item Datele cursului (SNP, SNG, SNN față de Brent, 2015--2026): corelația medie 0,13, max 0,54 (2020), min $-$0,32 (2022) --- aceeași cronologie.
  \item Decizia: în afara eșantionului de estimare (după 12.2020), reducerea varianței unui coș de acțiuni energetice românești acoperit cu Brent este de $-$1,7\,\% (acoperire dinamică) și de $-$6,9\,\% (statică; o valoare negativă înseamnă risc mai mare) --- riscul acestora este de reglementare și de piață locală, nu legat de petrol.
\end{itemize}
